L'espace est rapporté à un repère orthonormé \Oijk.

On considère :

\begin{itemize}
\item $\alpha$ un réel quelconque ;
\item les points A(1~;~1~;~0), B(2~;~1~;~0) et C$(\alpha~;~3~;~\alpha)$ ;
\item $(d)$ la droite dont une représentation paramétrique est :

\[
\left\{
\begin{array}{l}
x = 1 + t\\
y = \phantom{11} 2t\\
z = \phantom{1}- t
\end{array}
\right., \quad t \in \mathbb{R}
\]
\end{itemize}

Pour chacune des affirmations suivantes, préciser si elle est vraie ou fausse, puis justifier la réponse donnée. Une réponse non argumentée ne sera pas prise en compte.

\textbf{Affirmation 1 :} Pour toutes les valeurs de $\alpha$, les points $A$, $B$ et $C$ définissent un plan et un vecteur normal à ce plan est $\vect{\jmath}\begin{pmatrix}0\\1\\0\end{pmatrix}$.

\textbf{Affirmation 2 :} Il existe exactement une valeur de $\alpha$ telle que les droites $(AC)$ et $(d)$ soient parallèles.

\textbf{Affirmation 3 :} Une mesure de l'angle $\widehat{\text{OAB}}$ est $135\degres$.

\textbf{Affirmation 4 :} Le projeté orthogonal du point $A$ sur la droite $(d)$ est le point H(1~;~2~;~2).

\textbf{Affirmation 5 :} La sphère de centre O et de rayon 1 rencontre la droite $(d)$ en deux points distincts.

On rappelle que la sphère de centre $\Omega$
et de rayon $r$ est l'ensemble des points de l'espace situés à une distance $r$ de $\Omega$.
